\begingroup
\selectlanguage{english}
\noindent\textbf{Paths, comparison, and the information carried by a model.}

These pages bring the long construction of categorical paths to its decisive
test.  A path should do more than give a word in arrows and their inverses.  Its
category should have the homotopy type of a space of paths, should vary with
both endpoints, and should turn an arbitrary functor into a map whose strict
fibres compute homotopy fibres.  Grothendieck wants these properties before
making any fibrant replacement of the category in which the paths lie.  The
problem is therefore one of comparison between concrete models, with their
actual arrows and their actual endpoint maps
\citep[ms.~108--171]{GrothendieckDerivateursVII}.

There are three models in play.  The category $\Ch(X)$ retains a finite interval
type as part of a path, and permits changes of type.  Concatenation is strictly
associative there.  A chosen sequence of longer alternating intervals gives
categories $C_n$ and a filtered colimit $C_\infty$; identity arrows added at the
ends become identifications in this colimit.  Between them stands the
Grothendieck construction $C_{\mathbf N^o}$, which retains the stage and the
transition morphisms.  Passing the diagram $(C_n)$ through homotopy first gives
yet another object, the ind-object $H_T$.  None of these operations should be
silently substituted for another.  The opening pages, from manuscript page~108, are largely
an investigation of what is lost by doing so.

For a filtered diagram the comparison from its Grothendieck construction to
its ordinary colimit is a weak equivalence.  Nerves commute with filtered
colimits, and filtered homotopy colimits of simplicial sets map by weak
equivalences to ordinary colimits.  This explains one side of the comparison
span on manuscript page~108.  The other side depends on the indexing category
and on the diagram.  Grothendieck correctly observes that cofinality in the
preorder of objects does not make the original indexing category contractible:
a one-object group already distinguishes the two statements.  The appropriate
general machinery is the homotopy-colimit theorem for the Grothendieck
construction and the comma-category criteria for homotopy cofinality
\citep{Thomason1979,Quillen1973}.

The proposed independence of the chosen sequence is more delicate than the
text acknowledges.  An isomorphism of ind-objects in the ordinary homotopy
category does not control the compatibility of the homotopies used at
successive stages.  Our finite-injection example makes this precise within the
general hypotheses of manuscript pages~99--101.  To a finite set we associate
the one-object category of its permutation group.  Injections act by extending
permutations by the identity.  Parallel injections induce naturally isomorphic
functors, yet a morphism between two cofinal sequences induces, on the colimit,
the non-surjective shift of the finitary permutation group.  Thus the general
comparison assertion on pages~117--121 is false.  This does not invalidate the
geometric path comparison proved later by a different argument.

The same distinction clarifies the questions about
$\Hot\longrightarrow\Hot_0^{\wedge}$.  In this notation the subscript $0$
refers to a chosen category of finite tests; it is not a truncation to
homotopy $0$-types.  Restricting unbased representable functors to finite tests
is neither generally faithful nor generally conservative.  The locally inner
permutation-group example in the commentary on Chapter~V, manuscript page~52,
already treats conservativity.  Here a telescope of degree-two maps of circles supplies an
explicit failure of faithfulness: an essential map to $K(\mathbf Z,2)$ is
invisible on every finite stage.  Its obstruction is calculated in the
commentary, rather than attributed vaguely to a failure of coherence.  The
cellular and representability tools are classical
\citep[\S3F and Theorem~4.57]{Hatcher2002}.  In addition, a category with
finitely many arrows need not have a homotopy-finite nerve.  Compactness of a
test in the homotopy category requires its own hypothesis.  These qualifications
matter whenever a displayed bijection is used to infer an isomorphism or an
equality of maps.

\medskip\noindent\textbf{External Hom and the role of the parameter category.}

Section~11 returns to the reason that paths entered the theory of derivators.
For a relative category $(\mathcal M,\Sigma)$ one can always form the
prederivator
\[
 A\longmapsto\mathcal M^A[\Sigma_A^{-1}],
\]
where $\Sigma_A$ is defined objectwise.  Grothendieck asks when the morphisms
between constant diagrams are represented by a homotopy type, with the expected
dependence on the homotopy type of $A$.  This is a substantive condition on a
relative category.  It does not follow merely from the existence of an
ordinary localization, and the existence of a prederivator is not the
existence of all the Kan extensions required of a derivator.

The discussion of a categorical two-sphere on pages~132--135 identifies a real
problem but gives an incorrect explanation.  A fixed finite circle-shaped
parameter can already support arbitrarily high degrees when the length of the
target path is allowed to increase: concatenate the same family of loops
repeatedly.  We work this out for the six-object suspension of the four-object
circle poset.  The same example gives a genuine failure of homotopy invariance
by comparing that circle poset with the one-object category $B\mathbf Z$.
Every functor from a group to a poset is constant; the resulting localization
has lost the family that the circle-shaped parameter sees.  The example
separates the valid concern about parameter categories from the proposed
bound on degrees.

There is a second distinction at page~135.  The total arrow category
$\operatorname{Fun}([1],C)$ has the homotopy type of $C$, for every $C$.
Evaluation and the constant-diagram functor give the needed adjunction and
natural transformation.  What fails for an arbitrary category is the claim
that its strict fixed-endpoint arrow sets are derived path spaces.  Likewise,
Thomason fibrancy of the target alone does not turn every ordinary functor
category into a derived internal Hom.  The commentary supplies a fibrant-poset
counterexample and distinguishes a derived mapping-space construction from
this incorrect assertion.  The separate claim that ordinary functor categories
preserve Thomason fibrancy is also false.  An explicit equivariant horn-cell
construction gives a fibrant category $E$ for which
$\operatorname{Fun}(B(\mathbf Z/2),E)$ has a nonfibrant retract; the proof uses
the fixed-point pushout formulas of
\citet[Lemmas~4.7--4.8]{MayZakharevichStephan2017}.
Simplicial localization provides a systematic
external Hom for a relative category, while additional model-theoretic
hypotheses justify comparisons with categories of ordinary diagrams
\citep{DwyerKan1980Simplicial,DwyerKan1980Calculating,Groth2013,Cisinski2009LocallyConstant,Lurie2009HTT}.

\medskip\noindent\textbf{What the path comparison establishes.}

Sections~12 and~13 concern $\Sigma=\operatorname{Mor}(X)$.  Alternating paths
of unbounded length are the essential construction.  By choosing the parity
of their end arrows one obtains finite-stage Grothendieck fibrations or
opfibrations over the endpoint category.  The categories of lifts pass to
filtered unions.  This proves smoothness and properness of
$\Ch_\infty(X)\to X\times X$, although the union need not itself be a
Grothendieck fibration or opfibration.  The lift-category argument, not the
existence of compatible Cartesian arrows at all stages, is what survives the
passage to the union.  The terminology comes from base-change phenomena:
modern asphericity theory makes the categorical tests precise
\citep{Maltsiniotis2005Asphericite}.

For the category of all finite paths the source proves smoothness by a useful
explicit construction.  Given endpoint arrows and a path $c$, insert those
arrows at the two ends to form $\widetilde c$.  The category of lifts is
compared with the slice over $\widetilde c$, which has a terminal object.
Padding, retraction, and their naturality give the contraction.  This is a
sound argument once the notation in the diagrams is read correctly.  Its
claimed dual, however, does not prove properness.  An empty path admits no
outgoing change to a nonempty interval type.  Over $X=[1]$, changing only one
endpoint therefore gives an empty category of outgoing lifts.  The same
obstruction applies to $C_{\mathbf N^o}$.  The assertion that both these
endpoint projections are proper must remain false in the audit.

The desired comparison nevertheless survives.  For $\Ch(X)$, concatenating an
endpoint arrow, or its reversed path, defines transport between fibres.
The two composites are connected to the identity by the explicit
backtracking-cancellation zigzags.  Combining these equivalences with the
proved smoothness gives the hypotheses of the dual of Quillen's Theorem~B.
Hence strict endpoint fibres compute homotopy fibres.  The known total and
fibrewise equivalences
$C_{\mathbf N^o}\to\Ch_\infty(X)$ transfer the same property to the middle
model.  Finally the total equivalence
$C_{\mathbf N^o}\to\Ch(X)$ becomes a fibrewise equivalence by comparison of
homotopy fibres.  This repairs the concluding argument without assigning a
false properness property to either total category.

It follows that the two geometric path models have equivalent endpoint
fibres.  Strict concatenation in the category of finite paths then furnishes
the expected composition in $\Hot$, including the loop-space multiplication.
The diagonal factorization and the mapping-path construction on pages~170--171
work with the perfect model $\Ch_\infty(X)$; their formal proof is expanded in
the commentary.  Evrard's earlier categorical cocylinder and Shoikhet's
refinement give closely related factorization theorems, with precisely stated
homotopy-fibration conclusions.  They should not be cited as asserting that
Grothendieck's particular maps are Thomason fibrations or that the inclusion
of constant paths is a Thomason cofibration \citep{Evrard1975,Shoikhet2016}.

The larger programme is continuous with \emph{Pursuing Stacks}, \S69: homotopy
theory is to retain the diagram categories needed to form and compare fibres,
cofibres, and their operations.  The letter to Thomason of 2~April~1991
explicitly returns to unbounded categorical paths and to constructions that
avoid a preliminary Kan replacement
\citep{Grothendieck1983,Grothendieck1991Letter}.  Here the useful boundary is
precise.  The geometric comparison and mapping-path factorization can be
proved; the general finite-test and ordinary-internal-Hom assertions require
correction.  The developments marked \ours\ give the explicit counterexamples
and repairs used in the commentary.  They record arguments given here,
with the literature checks stated in their blocks, and make no claim of
priority.

\medskip\noindent\textit{Commentary, audit, and mathematical developments
AI-produced by GPT-6 Astra (OpenAI, ChatGPT Pro mode), checked and integrated
by Claude (Anthropic).}

\endgroup
